% Generado por scripts/tablas.py a partir de Practica3_pendulos_acoplados.xlsx. No editar a mano.

\begin{table}[ht]
\centering
\caption{Constante del muelle obtenida con cada método al repetir el ajuste sin las longitudes mayores. Con tres puntos el ajuste tiene un solo grado de libertad, así que estos valores solo indican tendencias.}
\label{tab:robustez}
\renewcommand{\arraystretch}{1.25}
\begin{tabular}{@{}l c c c@{}}
\toprule
Longitudes ajustadas & $k$ de $\omega_1^{2}$ (\si{\newton\per\meter}) & $k$ de $\omega_{-}$ (\si{\newton\per\meter}) & $k$ de $\omega_{+}$ (\si{\newton\per\meter}) \\
\midrule
$\ell_1$ a $\ell_5$ & \num{1.73} & \num{2.48} & \num{1.84} \\
$\ell_1$ a $\ell_4$ & \num{1.85} & \num{2.35} & \num{2.21} \\
$\ell_1$ a $\ell_3$ & \num{2.76} & \num{2.61} & \num{3.12} \\
\bottomrule
\end{tabular}
\end{table}
